\section{Phân tích overfitting / underfitting}\label{sec:analysis}

\paragraph{Cách đọc.} Các đường cong dưới đây vẽ loss (cross-entropy) và macro F1 theo epoch trên tập train (đã undersample) và
\code{val\_us} (validation undersample \emph{cùng tỉ lệ}), nên hai đường so sánh trực tiếp được. Chấm tròn là epoch được early
stopping giữ lại. Lưu ý: số liệu train là trung bình \emph{trong lúc huấn luyện} (dropout đang bật, trọng số thay đổi trong epoch),
nên với mô hình có dropout, F1 train có thể thấp hơn F1 validation. Bảng~\ref{tab:gap} tóm tắt các dấu hiệu chính.

\input{\tabdir/tab_gap}

\begin{figure}[H]\centering
\includegraphics[width=\linewidth]{curves_A_rnn.png}
\includegraphics[width=\linewidth]{curves_A_cnn.png}
\caption{Biểu diễn A: RNN (trên) và CNN đơn giản / + BatchNorm / + Dropout (dưới).}\label{fig:curves_A}
\end{figure}

\subsection{Biểu diễn A: giới hạn bởi thông tin, không phải bởi năng lực mô hình}
\begin{itemize}
\item Khoảng cách train--val rất nhỏ ($\le$0,02 macro F1; loss 0,20 so với 0,23) ở mọi cấu hình, và \textbf{tăng năng lực mô hình}
  (2 tầng, 128 units, hai chiều, 3 lớp conv, gấp đôi filters) \textbf{không cải thiện} validation. Đây không phải overfitting, cũng
  không phải underfitting do mô hình quá nhỏ: một dòng đơn lẻ không chứa đủ thông tin để tách ransomware khỏi white (sai số
  không giảm được --- \emph{irreducible error}). XGBoost trên cùng thông tin cũng chỉ đạt 0,575.
\item \textbf{Overfit nhẹ}: \code{A\_gru\_nodrop} và \code{A\_cnn\_bn} có F1 train cao nhất (0,77) và val loss đạt cực tiểu sớm hơn
  epoch tốt nhất (51 so với 63; 47 so với 53) rồi tăng lại (0,232 và 0,229 ở epoch cuối) trong khi train loss tiếp tục giảm --- dấu
  hiệu kinh điển của overfitting. Early stopping dừng kịp; trên test, mức overfit này gần như không gây hại.
\item \textbf{Underfit}:
  \begin{itemize}
  \item \code{A\_cnn\_do}, \code{A\_cnn\_bn\_do}, \code{A\_mlp}: F1 train \emph{thấp hơn} validation, cả hai đều thấp hơn bản không
    dropout. Dropout 0,3 với mô hình chỉ $\approx$23 nghìn tham số là điều chuẩn thừa.
  \item \code{A\_gru\_lr3e-4}: learning rate nhỏ $\Rightarrow$ đường cong tăng chậm, epoch tốt nhất là 84, cả train và val đều thấp
    hơn (0,71/0,71). \code{A\_gru\_bs2048}: batch lớn $\Rightarrow$ số bước cập nhật/epoch giảm 4 lần, epoch tốt nhất 98/100 ---
    \emph{chưa hội tụ} khi hết ngân sách epoch (Hình~\ref{fig:curves_reg}).
  \end{itemize}
\end{itemize}

\begin{figure}[H]\centering
\includegraphics[width=\linewidth]{curves_A_reg.png}
\includegraphics[width=\linewidth]{curves_A_lr.png}
\caption{Biểu diễn A: ảnh hưởng của dropout và độ phức tạp (trên); learning rate và batch size (dưới).}\label{fig:curves_reg}
\end{figure}

\subsection{Biểu diễn B: overfitting rõ hơn và vai trò của điều chuẩn}
\begin{itemize}
\item Với lịch sử, mô hình có nhiều thông tin hơn nên F1 train lên tới 0,91--0,92, và \textbf{khoảng cách train--val lớn hơn hẳn}
  (GRU: 0,911 so với 0,855; val loss 0,169 so với train 0,095 ở epoch cuối). Val loss của \code{B\_gru} đạt cực tiểu ở epoch 22 rồi
  tăng dần, trong khi val F1 vẫn đi ngang/tăng nhẹ đến epoch 53: mô hình tiếp tục \emph{tự tin quá mức} vào các dự đoán sai
  (loss tăng) dù số dự đoán đúng không giảm (F1 không giảm). Vì đề bài đánh giá bằng F1, early stopping theo F1 thay vì loss
  là lựa chọn hợp lý; trọng số lớp $w$ hiệu chỉnh sau cũng bù một phần cho xác suất bị lệch.
\item Nguyên nhân overfit dễ hơn: các dòng của cùng một địa chỉ có lịch sử chồng lấp nhau, nên số mẫu ``độc lập'' thực sự ít hơn
  nhiều so với 126 nghìn dòng (trong toàn bộ dữ liệu, 3.403 địa chỉ CryptoWall/CryptoLocker tạo ra 21.705 dòng, trong khi 9.223 dòng Cerber đến từ 9.177 địa chỉ); mô hình có thể học thuộc
  mẫu riêng của từng địa chỉ trong tập train.
\item \textbf{CNN không điều chuẩn} (\code{B\_cnn\_simple}) overfit nhanh nhất: epoch tốt nhất là 12, chênh F1 tăng từ 0,036 lên
  0,071 ở epoch cuối, test chỉ 0,686. \textbf{Thêm BatchNorm + Dropout} (\code{B\_cnn\_bn\_do}): train/val loss gần trùng nhau
  (0,131/0,136), chênh F1 0,016 --- tổng quát hoá tốt nhất --- và test tăng lên 0,724. Đây là điểm khác so với biểu diễn A:
  \emph{điều chuẩn chỉ có ích khi mô hình thực sự có xu hướng overfit}.
\item GRU 2 tầng (dropout giữa các tầng) có khoảng cách nhỏ hơn GRU 1 tầng (0,034 so với 0,056) và tốt hơn trên test --- dropout
  giữa tầng đóng vai trò điều chuẩn. L=32 overfit sớm hơn (dừng ở epoch 26).
\item \code{B\_mlp} (chỉ bước hiện tại) gần như không overfit nhưng cả train và val đều thấp (0,87/0,85): \emph{underfit} do
  thiếu thông tin lịch sử, không phải do mô hình nhỏ.
\end{itemize}

\begin{figure}[H]\centering
\includegraphics[width=\linewidth]{curves_B.png}
\includegraphics[width=\linewidth]{curves_B_len.png}
\caption{Biểu diễn B: RNN vs CNN vs MLP chỉ dùng bước hiện tại (trên); GRU với độ dài lịch sử L = 1, 4, 16, 32 (dưới).}
\label{fig:curves_B}
\end{figure}

\subsection{Biện pháp đã áp dụng và đề xuất}
\begin{itemize}
\item \textbf{Đã áp dụng}: early stopping (patience 10, khôi phục trọng số tốt nhất) cho mọi mô hình; Dropout ở đầu phân loại và
  giữa các tầng RNN; BatchNorm và Dropout cho CNN; cắt gradient; undersampling + hiệu chỉnh trọng số lớp trên validation.
\item \textbf{Với overfit (biểu diễn B)}: weight decay / L2 (chưa thử --- \code{weight\_decay=0}); giảm số lần xuất hiện được dùng
  của các địa chỉ có rất nhiều dòng (để mỗi địa chỉ đóng góp cân bằng hơn); làm nhiễu lịch sử (bỏ ngẫu nhiên một số bước --- một
  dạng tăng cường dữ liệu cho chuỗi); label smoothing để giảm hiện tượng tự tin quá mức.
\item \textbf{Với underfit}: tăng learning rate hoặc dùng lịch giảm learning rate (cosine/one-cycle) thay vì lr nhỏ cố định; khi dùng
  batch lớn thì tăng lr tương ứng hoặc tăng ngân sách epoch; giảm dropout cho mô hình nhỏ. Với biểu diễn A, cách duy nhất để vượt giới
  hạn là \emph{thêm thông tin} (như biểu diễn B), không phải mô hình lớn hơn.
\end{itemize}
